% Fixture: imágenes PNG (raster) y PDF (vectorial) en una subcarpeta,
% resueltas con \graphicspath y sin extensión explícita.
\documentclass[11pt,a4paper]{article}
\usepackage[T1]{fontenc}
\usepackage[utf8]{inputenc}
\usepackage[spanish]{babel}
\usepackage{graphicx}
\graphicspath{{figuras/}}

\title{Documento con imágenes}
\author{Fixture de pruebas}
\date{}

\begin{document}
\maketitle

La figura~\ref{fig:raster} es un PNG y la figura~\ref{fig:vector} es un PDF.

\begin{figure}[h]
  \centering
  \includegraphics[width=0.5\textwidth]{degradado}
  \caption{Imagen raster (PNG).}\label{fig:raster}
\end{figure}

\begin{figure}[h]
  \centering
  \includegraphics[width=0.4\textwidth]{diagrama}
  \caption{Imagen vectorial (PDF).}\label{fig:vector}
\end{figure}

\end{document}
